\section{Consumption and the permanent-income hypothesis}

\subsection*{(a) Testing the additional consumption lags}
Construct quarterly real consumption per capita as
\[
C_t=\frac{\text{real nondurables}_t+\text{real services}_t}
{\text{population}_t}.
\]
The two consumption components are already real, so no further CPI
deflation is applied. The data contain 220 quarters, 1948Q1--2002Q4.
Losing four initial observations leaves $n=216$, from 1949Q1 to 2002Q4.
OLS estimation of
\[
C_t=\beta_0+\beta_1C_{t-1}+\beta_2C_{t-2}
+\beta_3C_{t-3}+\beta_4C_{t-4}+\varepsilon_t
\]
gives the estimates in Table~\ref{tab:ar}.
\begin{table}[H]
\centering
\caption{Consumption in levels: four autoregressive lags}
\label{tab:ar}
\begin{tabular}{lrrr}
\toprule
Regressor & Estimate & Conventional SE & White HC0 SE \\
\midrule
Constant & 0.00002439 & 0.00001925 & 0.00001940 \\
$C_{t-1}$ & 1.192335 & 0.067378 & 0.067873 \\
$C_{t-2}$ & -0.128134 & 0.105701 & 0.102071 \\
$C_{t-3}$ & 0.144329 & 0.105489 & 0.107384 \\
$C_{t-4}$ & -0.207464 & 0.067530 & 0.073593 \\
\midrule
Observations & 216 & & \\
\bottomrule
\end{tabular}

\end{table}
For $H_0:\beta_2=\beta_3=\beta_4=0$, the usual restricted-versus-unrestricted
test is
\[
F=\frac{(SSR_R-SSR_U)/3}{SSR_U/(216-5)}
=8.409997,\qquad F_{3,211}\text{ p-value}=0.00002636.
\]
We reject at $5\%$: the additional lags have joint predictive power,
contrary to the restriction implied by Hall's hypothesis in this model.
The heteroskedasticity-robust HC0 Wald test agrees:
$\chi^2_3=20.670589$, with $p=0.00012323$.

\subsection*{(b) Why current income growth is endogenous}
Let $\Delta C_t=C_t-C_{t-1}$ and $\Delta Y_t=Y_t-Y_{t-1}$.
In
\[
\Delta C_t=\gamma_0+\gamma_1\Delta Y_t+\varepsilon_t,
\]
the error includes news about permanent income. Such news can also affect
current income, so generally
$\operatorname{Cov}(\Delta Y_t,\varepsilon_t)\ne0$.
With a constant and suitable laws of large numbers,
\[
\operatorname{plim}\widehat\gamma_{1,OLS}
=\gamma_1+
\frac{\operatorname{Cov}(\Delta Y_t,\varepsilon_t)}
{\operatorname{Var}(\Delta Y_t)}.
\]
Thus OLS is generally inconsistent. The direction of the bias depends
on the covariance; it does not follow from the model alone.

\subsection*{(c) Two-stage least squares in log changes}
Define
\[
c_t=\log C_t,\qquad
y_t=\log\!\left(
\frac{\text{real disposable income}_t}{\text{population}_t}
\right).
\]
Estimate $\Delta c_t=\gamma_0+\gamma_1\Delta y_t+\varepsilon_t$
using the instrument vector
\[
z_t=(1,\Delta c_{t-2},\Delta c_{t-3},
             \Delta c_{t-4},\Delta c_{t-5})'.
\]
For example, $\Delta c_{t-2}=\log(C_{t-2}/C_{t-3})$,
so the oldest required level is $C_{t-6}$.
The resulting sample has $n=214$ observations, 1949Q3--2002Q4.

Let $X$ contain a constant and $\Delta y_t$, let $Z$ contain the five
instruments, and let $\widehat X=P_ZX$.
The estimator and White HC0 covariance matrix are
\[
\widehat\gamma=(X'P_ZX)^{-1}X'P_Z\Delta c,\qquad
\widehat{\operatorname{Var}}_{HC0}(\widehat\gamma)
=(\widehat X'\widehat X)^{-1}
\widehat X'\operatorname{diag}(\widehat\varepsilon_t^2)\widehat X
(\widehat X'\widehat X)^{-1}.
\]
Here the residuals are the structural residuals
$\widehat\varepsilon=\Delta c-X\widehat\gamma$.
No degrees-of-freedom correction is applied.
\begin{table}[H]
\centering
\caption{Consumption growth on current income growth, $n=214$}
\label{tab:iv}
\begin{tabular}{lrrrr}
\toprule
 & \multicolumn{2}{c}{OLS} & \multicolumn{2}{c}{2SLS} \\
\cmidrule(lr){2-3}\cmidrule(lr){4-5}
Regressor & Estimate & White HC0 SE & Estimate & White HC0 SE \\
\midrule
Constant & 0.003653 & 0.000397 & 0.002613 & 0.001157 \\
$\Delta\log Y_t$ & 0.239665 & 0.041784 & 0.436925 & 0.202555 \\
\midrule
Observations & 214 & & 214 & \\
\bottomrule
\end{tabular}

\end{table}
The IV slope is $0.436925$ (White SE $0.202555$).
A one-percentage-point increase in income growth is associated with
approximately $0.437$ percentage points more consumption growth.
The usual asymptotic two-sided test of a zero slope gives $p=0.0310$.
The $95\%$ normal confidence interval is $[0.0399,0.8339]$.
The estimate is consistent with an economically relevant response
to current income. In the stylized Campbell--Mankiw interpretation,
it suggests a substantial rule-of-thumb component; identifying the
log-regression coefficient literally with a population fraction requires
additional aggregation and scaling assumptions.

Instrument relevance is limited: the conventional first-stage test
of the four excluded instruments gives $F_{4,209}=2.509115$
(first-stage $R^2=0.045821$).
The HC0 excluded-instrument Wald statistic is $\chi^2_4=12.832405$.
This weak first stage makes the conventional IV normal approximation
and the specification tests below less reliable and potentially low-powered.

\clearpage
\subsection*{(d) Exogeneity and overidentifying restrictions}
For exogeneity, first regress $\Delta y_t$ on $z_t$ and obtain
$\widehat v_t$. Then estimate
\[
\Delta c_t=\gamma_0+\gamma_1\Delta y_t+\rho\widehat v_t+u_t.
\]
Under valid instruments, testing $\rho=0$ is a regression-form
Durbin--Wu--Hausman test of income-growth exogeneity.
The HC0 estimate is $\widehat\rho=-0.206732$ with SE $0.193630$.
The robust Wald statistic is
\[
W=\frac{\widehat\rho^2}{\widehat{\operatorname{Var}}_{HC0}(\widehat\rho)}
=1.139904,\qquad p=0.285673\quad(\chi^2_1).
\]
We do not reject exogeneity at $5\%$.

For overidentification, use a heteroskedasticity-robust Hansen test.
Starting with the 2SLS residuals from (c), form
\[
\widehat S=\frac1n\sum_t z_tz_t'\widehat\varepsilon_t^2,\qquad
\bar g(\gamma)=\frac1n Z'(\Delta c-X\gamma).
\]
Re-estimate by minimizing
\[
\bar g(\gamma)'\widehat S^{-1}\bar g(\gamma).
\]
The two-step efficient GMM estimate is
$\widehat\gamma_{GMM}=(0.00309653,0.35328490)'$.
The minimized criterion gives
\[
J=n\bar g(\widehat\gamma_{GMM})'\widehat S^{-1}
\bar g(\widehat\gamma_{GMM})
=5.087017,\qquad p=0.165535\quad(\chi^2_{5-2}).
\]
We do not reject the three overidentifying restrictions.
Re-estimation matters: evaluating this robust criterion at an
inefficient 2SLS estimate does not generally yield the usual
$\chi^2_3$ Hansen test.
\begin{table}[H]
\centering
\caption{Main specification tests}
\label{tab:tests}
\begin{tabular}{lrrr}
\toprule
Test & Statistic & df & $p$-value \\
\midrule
Hall lag restrictions ($F$) & 8.409997 & $(3,211)$ & 0.00002636 \\
Exogeneity (HC0 $\chi^2$) & 1.139904 & $1$ & 0.285673 \\
Overidentification (Hansen $J$) & 5.087017 & $3$ & 0.165535 \\
\bottomrule
\end{tabular}

\end{table}
Failure to reject does not establish exogeneity or instrument validity,
particularly with the weak first stage.
White's covariance estimator follows the requested method; its
large-sample calibration assumes the moment errors have no remaining
serial correlation. If that condition fails, a HAC covariance estimator
would be required.
